% Fragmento integrable. Preambulo anfitrion: amsmath, amssymb, longtable, hyperref.
\section{Isolement de la limite et changement d’échelle du sélecteur global}
\label{hmtfg:fragmento:aislamiento-del-limite-y-escalado-del-selector-global}

\subsection{Continuation indéfinie de la première racine et exclusion de l’intervalle intermédiaire}
\label{hmtfg:escalado:15}

La classification complète démontrée précédemment conserve la famille et le sélecteur de première racine nouvelle. Le transport ordonné conserve les zéros et leur minimum ; le mot de retours conserve la lecture chronologique nonadique.

\subsubsection{Classification et existence à toute profondeur}
\label{hmtfg:escalado:15.1}

Écrivons \(W_n=\tau_1\cdots\tau_{2^n-1}\), avec \(\tau_j=(-1)^{v_2(j)}\), et \(a_*=\min\Lambda_\tau\). L’existence et la compacité de \(\Lambda_\tau\) sont démontrées dans le développement de prolongation, \ref{hmtfg:prolongacion:6.1}–\ref{hmtfg:prolongacion:6.2} ; son minimum existe par compacité et ordre.

\textbf{Classification.} Toute application strictement unimodale à maximum, sur un intervalle invariant, dont le point critique a pour période exacte \(2^n\) et dont l’itinéraire \(K\) satisfait \(K<\tau\), a pour itinéraire \(W_n0\).

L’induction s’effectue au moyen du retour effectif. Pour une période d’au moins huit, la comparaison avec \(\tau\), jointe aux intervalles qui piégeraient des orbites de signes incompatibles, impose le préfixe \((+,-,+,+,+,-)\). On en déduit
\[
J_1=[d_2,d_4],\quad F(J_1)=[d_3,1],\quad
d_4<d_3,\quad F^2(J_1)=J_1.
\]
Le retour normalisé par \(h(x)=d_2x\) a pour maximum un, pour période la moitié de la période initiale et pour itinéraire \(K'_j=-K_{2j}<\tau\). L’ordre est conservé parce que
\[
O_{2q-1}(K)=-O_{q-1}(K'),\qquad
K_{2q}-\tau_{2q}=-(K'_q-\tau_q).
\]
Les cas initiaux sont \(+0\) et \(+-+0\) ; l’induction reconstruit \(W_n0\). Pour les applications paires, on a aussi \(0<d_4<-d_2<1\), de sorte que le retour s’étend à \([-1,1]\) et coïncide avec \eqref{hmtfg:escalado:eq:6}.

Avant \(a_*\), tous les itinéraires sont inférieurs à \(\tau\) : les ensembles de comparaison stricte sont ouverts par la séparation démontrée dans \ref{hmtfg:prolongacion:6.1} ; l’intervalle précédant \(a_*\) est connexe et commence par \(+^\infty<\tau\).

\textbf{Existence de chaque première racine.} Supposons \(\lambda_{n-1}<a_*\) construite, et soit \(p=2^{n-1}\). Les zéros de \(c_p\) dans \([\lambda_{n-1},a_*]\) forment un ensemble fini non vide : la fonction est analytique et \(c_p(a_*)\ne0\). Soit \(z<a_*\) son dernier zéro. Dans \((z,a_*]\), \(c_p\) est de signe \(\tau_p\). Pour \(g_\lambda=f_\lambda^p\),
\[
g_\lambda'(0)=0,\qquad
c_{2p}(\lambda)=c_p(\lambda)+O(c_p(\lambda)^2).
\]
Au voisinage de \(z\), à droite, \(c_{2p}\) est de signe \(\tau_p\) ; en \(a_*\), son signe est opposé. Il existe un zéro de \(c_{2p}\) dans \((z,a_*)\), et sa période est exactement \(2p\), car on y a \(c_p\ne0\). La finitude des zéros fournit la première racine nouvelle dans \((\lambda_{n-1},a_*)\). Le dernier zéro auxiliaire prouve l’existence ; le choix effectif reste la première racine. Le cas de base est \(\lambda_1=1/u_0(1)<a_*\).

\subsubsection{Limite de la suite initiale}
\label{hmtfg:escalado:15.2}

\textbf{Théorème.} Le sélecteur initial existe à tous les niveaux et satisfait
\begin{equation}
\boxed{K_{\lambda_n}=W_n0,\qquad
\lambda_n\nearrow a_*=\min\Lambda_\tau.}
\label{hmtfg:escalado:eq:50}
\end{equation}
\textbf{Démonstration.} Le résultat d’existence précédent donne une suite croissante majorée par \(a_*\). La classification fixe tous ses itinéraires. Soit \(b\le a_*\) sa limite. Pour chaque \(p\), les termes suffisamment avancés conservent les signes opposés \(\tau_p,\tau_{2p}\). La formule de Taylor et la borne uniforme \(B_p=16^p(16^p-1)/15\) donnent
\[
|c_{2p}(\lambda_n)-c_p(\lambda_n)|
\le\tfrac12 B_p|c_p(\lambda_n)|^2,\qquad
|c_p(\lambda_n)|>2/B_p.
\]
Le passage à la limite conserve \(\operatorname{sign}c_p(b)=\tau_p\), avec une valeur absolue d’au moins \(2/B_p>0\). Pour tout \(p\), cela implique \(b\in\Lambda_\tau\) ; la minimalité donne \(b=a_*\).

Le résultat établit l’existence indéfinie, la conservation de la cascade symbolique et la sélection de son premier paramètre limite pour la règle globale initiale. La complétion HMT transporte cette suite et sa limite en conservant l’ordre et la lecture.

\subsubsection{Certificat d’exclusion de l’intervalle intermédiaire}
\label{hmtfg:escalado:15.3}

Le programme de pont certifie
\begin{equation}
\Lambda_\tau\cap[\lambda_{8,\rm inf},1.874038]=\varnothing,
\label{hmtfg:escalado:eq:51}
\end{equation}
où
\[
\lambda_{8,\rm inf}
=1.874027744154450672897007573595092408435133414523352839687903472138975
\]
est l’extrémité inférieure rationnelle de l’intervalle certifié de \(\lambda_8\).

Le recouvrement contient 374 intervalles exactement adjacents et aucune région en suspens : 319 avec \(c_{512}>0\), 10 avec \(c_{1024}<0\), deux avec \(c_{2048}>0\), tous opposés au signe de \(\tau\), et 43 proches de centres de périodes 256, 512 ou 1024, exclus par leur deuxième retour.

Dans ce dernier cas, \(g=f_\lambda^p\), \(d=g(0)\) et une borne uniforme \(|g''|\le M\) entre zéro et \(d\) satisfont \(M|d|<2\). Alors
\[
|g(d)-d|\le\tfrac12M|d|^2<|d|\quad(d\ne0).
\]
Les deux retours ont le même signe ; pour \(d=0\), ils sont tous deux nuls. Les deux situations excluent \(\tau_{2p}=-\tau_p\) avec des signes stricts.

L’évaluation utilise des moments rationnels, des queues géométriques du potentiel et de sa dérivée seconde, et \(|u_0'''|<6\). Le centrage paramétrique intersecte l’enveloppe naturelle avec
\[
c_j(\lambda_{\rm medio})+
\partial_\lambda c_j(I)\,(I-\lambda_{\rm medio});
\]
la dérivée est propagée sur la boîte entière. La propriété d’autoapplication de \([-1,1]\) est démontrée avant d’y restreindre les enveloppes.

Le certificat complet conserve le recouvrement et vérifie son pavage exact, ainsi que chaque inégalité de signe ou de Taylor.

\subsubsection{Localisation conjointe}
\label{hmtfg:escalado:15.4}

Le croisement local \(\lambda_*\) réalise \(\tau\) par ses retours réels à toute profondeur, de sorte que \(a_*\le\lambda_*\). Par \eqref{hmtfg:escalado:eq:50}, \(\lambda_8<a_*\) ; avec \eqref{hmtfg:escalado:eq:51},
\begin{equation}
\boxed{1.874038<a_*\le\lambda_*<1.874039.}
\label{hmtfg:escalado:eq:52}
\end{equation}
La localisation conserve l’itinéraire canonique et situe son premier paramètre limite dans le même intervalle que le croisement stable.

\subsection{Isolement par première sortie et coïncidence des limites}
\label{hmtfg:escalado:16}

Cette section complète l’isolement identifié dans la section~\ref{hmtfg:escalado:15.4}. Elle conserve la famille HMT \eqref{hmtfg:escalado:eq:1}, la sélection minimale \eqref{hmtfg:escalado:eq:50}, la carte \eqref{hmtfg:escalado:eq:5} et la feuille stable des sections~\ref{hmtfg:escalado:5}--\ref{hmtfg:escalado:7}. La preuve utilise l’ordre \(a_*\le\lambda_*\) : il suffit d’exclure une séparation vers le côté négatif du cône expansif.

\subsubsection{Bornes supérieure et inférieure du transport transversal}
\label{hmtfg:escalado:16.1}

Outre \eqref{hmtfg:escalado:eq:9}, le lemme 2.1 de \href{https://arxiv.org/pdf/1410.3277}{Hertling–Spandl, section 2} fournit \(\|D\Phi\|<0.9\) dans la même boule de rayon \(0.01\), où
\[
\Phi_u(u,\nu)=u-\frac{U(u,\nu)-u}{3.669}.
\]
Par conséquent,
\begin{equation}
\left|1-\frac{U_u-1}{3.669}\right|<0.9,\qquad
U_u<1+1.9(3.669)=7.9711.
\label{hmtfg:escalado:eq:53}
\end{equation}
Le nombre \(3.669\) appartient au préconditionneur publié de ce certificat analytique ; la famille générée \eqref{hmtfg:escalado:eq:1} conserve ses coefficients antérieurs.

Prenons \(\kappa=0.21\). Pour deux points de la boule joints par une sécante avec
\(\Delta u<0\), \(\|\Delta\nu\|_1\le\kappa|\Delta u|\), l’intégration des blocs de \(DR\) sur le segment donne
\begin{equation}
4.4034|\Delta u|
\le|\Delta u'|\le8.0887|\Delta u|,\qquad \Delta u'<0,
\label{hmtfg:escalado:eq:54}
\end{equation}
\begin{equation}
\|\Delta\nu'\|_1
\le(0.756+0.719\kappa)|\Delta u|
=0.90699|\Delta u|
<\kappa(4.4034)|\Delta u|.
\label{hmtfg:escalado:eq:55}
\end{equation}
Ici \(4.4034=4.521-0.560\kappa\) et \(8.0887=7.9711+0.560\kappa\). Ainsi, le cône des sécantes conserve son orientation et se dilate tant que le segment reste dans la boule.

\subsubsection{Réduction d’une séparation infinie à une région finie}
\label{hmtfg:escalado:16.2}

Supposons \(a_*<\lambda_*\), et posons
\[
f_n=R^{4+n}f_{a_*},\qquad
s_n=R^{4+n}f_{\lambda_*},\qquad
(\Delta u_n,\Delta\nu_n)=f_n-s_n.
\]
La borne positive de \(u'\) et \(\|\nu'\|_1/u'<0.144386<\kappa\) sur \(\Gamma(J)\) placent la sécante initiale dans le cône négatif. La feuille stable est écrite par rapport à \(g\), avec
\begin{equation}
s_n=g+e_n,\quad
\|e_n\|_{\mathcal A}<0.001666(0.83996)^n,\quad
|(e_n)_u|\le0.16\|(e_n)_\nu\|_1.
\label{hmtfg:escalado:eq:56}
\end{equation}
Les contrôles d’entrée impliquent
\begin{equation}
|\Delta u_0|
<0.004993128+0.00004
 +0.16(0.001396077+0.00004)
=0.00526290032< h,\qquad h=0.006.
\label{hmtfg:escalado:eq:57}
\end{equation}
Tant que \(|\Delta u_n|<h\), \eqref{hmtfg:escalado:eq:55}–\eqref{hmtfg:escalado:eq:56} donnent
\begin{equation}
\|f_n-g_0\|_{\mathcal A}
<(1+\kappa)h+0.001666+0.00004
=0.008966<0.01.
\label{hmtfg:escalado:eq:58}
\end{equation}
Le point \(s_n\) et le segment entre les deux restent également dans cette boule convexe. L’expansion inférieure \eqref{hmtfg:escalado:eq:54} impose une première sortie finie \(N\), et la borne supérieure appliquée à l’étape précédente fournit
\begin{equation}
0.006\le-\Delta u_N<0.0485322,\qquad
\|\Delta\nu_N\|_1\le0.21|\Delta u_N|.
\label{hmtfg:escalado:eq:59}
\end{equation}
On utilise l’inégalité \(8.0887h=0.0485322\) ; le saut complet, et non seulement la surface de première sortie, est inclus.

Par conséquent, \(f_N\) appartient à la région suivante :
\begin{equation}
\begin{split}
\mathcal C_-=\{\,g+(d_u,d_\nu)+e:\;&
-0.0485322\le d_u\le-0.006,\\
&\|d_\nu\|_1\le0.21|d_u|,\quad
|e_u|+\|e_\nu\|_1\le0.001666,\\
&|e_u|\le0.16\|e_\nu\|_1\,\}.
\end{split}
\label{hmtfg:escalado:eq:60}
\end{equation}
Puisque \(a_*\) réalise \(\tau\), tous ses retours normalisés sont des autoapplications unimodales de \([-1,1]\) ayant le même itinéraire. Cette propriété découle de l’induction des retours, indépendamment de la distance de \(f_N\) à \(g_0\). C’est pourquoi le certificat suivant traite \(\mathcal C_-\) intersectée avec cette classe d’autoapplications.

\subsubsection{Exclusion certifiée de la région de première sortie}
\label{hmtfg:escalado:16.3}

Le certificat rationnel d’exclusion par première sortie, reproduit dans le supplément computationnel, couvre exactement l’intervalle de \eqref{hmtfg:escalado:eq:59} par 64 bandes rationnelles adjacentes. Toutes sont exclues ; aucune subdivision supplémentaire ni région en suspens n’intervient dans le résultat.

Parmi les 64 bandes, 54 présentent un signe contraire et dix un retour contractant de période 16 ou 32. La plus grande borne de Lipschitz de ces retours est inférieure à \(0.475146491\) ; l’horizon maximal d’exclusion est l’itérée 64.

L’évaluation conserve comme variables partagées les quarante premiers coefficients de la perturbation, ainsi qu’une borne explicite de leur queue infinie. Pour chaque bande, le centre est \(g_0\) déplacé dans la coordonnée \(u\). Si \(a,b\ge0\) bornent les sensibilités duales dans \(u,\nu\), le support de l’incertitude du point fixe et du résidu stable est
\begin{equation}
0.00004\max(a,b)
+0.001666\max\left(b,\frac{0.16a+b}{1.16}\right).
\label{hmtfg:escalado:eq:61}
\end{equation}
La deuxième expression s’obtient en maximisant \(a|e_u|+b\|e_\nu\|_1\) sous les deux contraintes de \eqref{hmtfg:escalado:eq:56}. On utilise ainsi l’orientation de la feuille stable démontrée précédemment, en plus de sa norme.

Soit \(x_j\) l’orbite du centre et \(L_j\) sa sensibilité linéaire à tous les coefficients partagés. Le modèle de Taylor prend la forme
\[
c_j=x_j+L_j(\Delta v)+r_j,\qquad |r_j|\le E_j.
\]
Si \(\rho_j\) borne \(|L_j(\Delta v)|+E_j\), la récurrence extérieure employée est
\begin{equation}
E_{j+1}\le |f_0'(x_j)|E_j
+\tfrac12\sup|f_0''|\rho_j^2
+\sup|\Delta f'|\rho_j.
\label{hmtfg:escalado:eq:62}
\end{equation}
Les dérivées sont bornées sur le segment des deux états. L’orbite centrale est vérifiée à l’intérieur de \([-1,1]\) ; les autres orbites appartiennent à cet intervalle par la propriété d’autoapplication incluse dans le domaine certifié. Les calculs utilisent un arrondi rationnel vers l’extérieur, y compris pour les queues géométriques.

Chaque bande est exclue par l’une de deux inégalités :

\par\noindent\textbf{1.}\enspace Une valeur critique a un signe strictement opposé à \(\tau_j\).
\par\noindent\textbf{2.}\enspace Pour \(p=2^k\), le retour \(H=f^p\) a pour constante de Lipschitz \(L_p<1\) sur le segment entre \(0\) et \(H(0)\). Alors
\[
|H(H(0))-H(0)|\le L_p|H(0)|.
\]
Si \(H(0)\ne0\), les deux valeurs ont le même signe ; si \(H(0)=0\), elles s’annulent toutes deux. Dans les deux cas, l’itinéraire strict \(\tau_{2p}=-\tau_p\) est exclu.

La constante \(L_p\) s’obtient en propageant le segment initial \([-r,r]\), \(r\ge|H(0)|\), avec les boîtes de l’orbite critique et en multipliant les bornes extérieures des dérivées. On borne les applications de la bande entière, en maintenant \eqref{hmtfg:escalado:eq:61}.

\subsubsection{Théorème de coïncidence du paramètre d’accumulation}
\label{hmtfg:escalado:16.4}

\textbf{Théorème.} Pour la famille uniforme HMT \eqref{hmtfg:escalado:eq:1}, la limite de la sélection initiale des premières racines coïncide avec le croisement stable certifié :
\begin{equation}
\boxed{\lambda_n\nearrow a_*=\lambda_*,
\qquad 1.874038<\lambda_*<1.874039.}
\label{hmtfg:escalado:eq:63}
\end{equation}

\textbf{Démonstration.} \eqref{hmtfg:escalado:eq:50}–\eqref{hmtfg:escalado:eq:52} donnent \(\lambda_n\uparrow a_*\le\lambda_*\), les deux paramètres étant dans \(J\). Si l’inégalité était stricte, \eqref{hmtfg:escalado:eq:54}–\eqref{hmtfg:escalado:eq:59} produiraient un retour \(f_N\in\mathcal C_-\) réalisant \(\tau\). Le recouvrement complet de la section~\ref{hmtfg:escalado:16.3} exclut précisément cette possibilité. Donc \(a_*=\lambda_*\).

La preuve contrôle toutes les profondeurs au moyen d’une première sortie finie et d’un recouvrement uniforme de sa région. L’identité \eqref{hmtfg:escalado:eq:63} relie la limite du sélecteur global au croisement transversal.


\subsection{Identification des centres initiaux et changement d’échelle global}
\label{hmtfg:escalado:17}

La dernière étape conserve la sélection de \eqref{hmtfg:escalado:eq:50}. La continuation locale et les paramètres initiaux sont comparés par leur période exacte et leur classe hybride ; la coïncidence s’obtient par une unicité dans un voisinage paramétrique fixe.

\subsubsection{Réalisation analytique de degré deux et transversalité}
\label{hmtfg:escalado:17.1}

Le point fixe \(g\) de la section~\ref{hmtfg:escalado:3} est pair, analytique, à point critique quadratique et de combinatoire stationnaire de doublement. Le théorème 2.1 de \href{https://annals.math.princeton.edu/wp-content/uploads/annals-v164-n3-p01.pdf}{de Faria–de Melo–Pinto}, appliqué à cette unique combinatoire, donne un ensemble limite à un seul élément, avec une extension holomorphe propre de degré deux entre disques topologiques emboîtés. Son attraction sur l’intervalle identifie cet élément à \(g\), puisque \(Rg=g\).

Le lemme 3.2 de la même source étend un nombre fixe de retours à un voisinage analytique complet de \(g\), avec des images de degré deux et un module positif. La norme \eqref{hmtfg:escalado:eq:5} contrôle la norme uniforme sur un voisinage complexe suffisamment petit de \([-1,1]\), car on y a \(|(x^2-1)/2.5|<1\). La convergence \eqref{hmtfg:escalado:eq:3} entre dans ce voisinage ; la dépendance analytique des retours fournit, pour un certain entier fixe \(d\), une famille
\begin{equation}
\mathcal F_\lambda=R^d f_\lambda
\label{hmtfg:escalado:eq:64}
\end{equation}
d’applications de degré deux, analytique dans un voisinage complexe de \(\lambda_*\). Les disques et le nombre de retours sont fixés avant de faire varier \(\lambda\) dans ce voisinage.

La tangente paramétrique appartient au cône expansif, tandis que la feuille stable satisfait la borne de pente de la section~\ref{hmtfg:escalado:5}. Les retours conservent cette séparation. La réalisation de degré deux identifie localement la feuille stable à la classe hybride de l’application infiniment renormalisable ; \eqref{hmtfg:escalado:eq:64} est donc transversale à cette classe.

Le changement d’espace analytique est justifié en détail dans \ref{hmtfg:clasificacion:8.2}, consacrée au transport de la transversalité. Une direction tangente à la classe hybride produirait des tangentes renormalisées décroissantes, par la convergence uniforme sur un disque hybride et l’estimation de Cauchy. La direction certifiée satisfait, en revanche, \(|\dot f(1)|=|\dot u|/10\) et croît sous les retours grâce au cône expansif. Cette évaluation commune des deux représentants prouve la transversalité sans supposer l’équivalence de leurs normes.

\subsubsection{Unicité dans un voisinage fixe}
\label{hmtfg:escalado:17.2}

Pour \(n>d\), la classification de la section~\ref{hmtfg:escalado:15} donne
\begin{equation}
K(\mathcal F_{\lambda_n})=W_{n-d}0.
\label{hmtfg:escalado:eq:65}
\end{equation}
Le redressement d’une application de degré deux conserve l’orbite critique ; \eqref{hmtfg:escalado:eq:65} identifie sa classe hybride à celle du centre quadratique canonique de période \(2^{n-d}\). Les bornes a priori de la suite réelle de doublement sont celles employées dans le théorème d’universalité de Lyubich.

Le \href{https://www.maths.tcd.ie/EMIS/journals/Annals/149_2/lyubich.pdf\#page=69}{théorème 7.4 de Lyubich, p. 387–388} donne une seule intersection avec chacune de ces classes, aux niveaux suffisamment élevés, sur la transversale analytique et dans un voisinage fixe du paramètre d’accumulation.

Par \eqref{hmtfg:escalado:eq:63}, \(\lambda_n\) entre finalement dans ce voisinage. La continuation locale de la section~\ref{hmtfg:escalado:8}, indexée par la même période, appartient à la même classe hybride. L’unicité implique
\begin{equation}
\boxed{\lambda_n=\widehat\mu_n\quad(n\ge n_0),}
\label{hmtfg:escalado:eq:66}
\end{equation}
où \(\widehat\mu_n\) désigne le centre local de période physique \(2^n\). Il s’agit d’une réindexation de \(\mu_j\) par un décalage fixe déterminé par les retours et la cible de la section~\ref{hmtfg:escalado:8}. La sélection initiale de \(\lambda_n\) reste intacte.

L’entier \(n_0\) est existentiel dans cette preuve. La continuation à tous les niveaux et le mot \(W_n0\) sont déjà démontrés dans \eqref{hmtfg:escalado:eq:50} ; \eqref{hmtfg:escalado:eq:66} apporte l’identité métrique nécessaire à la limite.

\subsubsection{Théorème de changement d’échelle pour les premières racines HMT}
\label{hmtfg:escalado:17.3}

\textbf{Théorème.} Pour la famille uniforme \eqref{hmtfg:escalado:eq:1}, soient \(\lambda_n\) les premières racines nouvelles successives de période critique exacte \(2^n\). Il existe \(C>0\), \(0<\theta<1\) et \(M<\infty\) tels que, pour tout \(n\) suffisamment grand,
\begin{equation}
\lambda_*-\lambda_n=C\delta_{\mathrm F}^{-n}(1+e_n),
\qquad |e_n|\le M\theta^n,
\label{hmtfg:escalado:eq:67}
\end{equation}
et
\begin{equation}
\boxed{\lim_{n\to\infty}
\frac{\lambda_{n-1}-\lambda_{n-2}}
{\lambda_n-\lambda_{n-1}}=\delta_{\mathrm F}.}
\label{hmtfg:escalado:eq:68}
\end{equation}

\textbf{Démonstration.} \eqref{hmtfg:escalado:eq:66} transporte le changement d’échelle local \eqref{hmtfg:escalado:eq:24} aux racines initiales ; le décalage fini des indices est absorbé dans \(C,M\). Son signe est positif parce que \(\lambda_n<\lambda_*\).

Le reste en puissance admet aussi une justification directe par l’holonomie des classes hybrides. Le lemme 7.3 de Lyubich donne une régularité \(C^{1+\beta}\), avec une dérivée non nulle, sur la transversale au point d’accumulation. La dynamique instable unidimensionnelle est analytique et a pour multiplicateur \(\delta_{\mathrm F}>1\) ; sa coordonnée linéarisante place les centres successifs à une distance proportionnelle à \(\delta_{\mathrm F}^{-n}\). La composition avec l’inverse de l’holonomie donne un terme linéaire non nul et un reste \(O(\delta_{\mathrm F}^{-(1+\beta)n})\). On peut ainsi prendre \(\theta=\delta_{\mathrm F}^{-\beta}\).

En écrivant
\[
r_n=\frac{\delta_{\mathrm F} e_{n-1}-e_n}{\delta_{\mathrm F}-1},
\qquad K=\frac{M(\delta_{\mathrm F}+\theta)}{\delta_{\mathrm F}-1},
\]
on obtient
\[
\lambda_n-\lambda_{n-1}
=C(\delta_{\mathrm F}-1)\delta_{\mathrm F}^{-n}(1+r_n),\qquad
|r_n|\le K\theta^{n-1}.
\]
Par conséquent,
\begin{equation}
\left|
\frac{\lambda_{n-1}-\lambda_{n-2}}
{\lambda_n-\lambda_{n-1}}-\delta_{\mathrm F}
\right|
\le
\frac{\delta_{\mathrm F} K(1+\theta)\theta^{n-2}}
{1-K\theta^{n-1}},
\quad K\theta^{n-1}<1.
\label{hmtfg:escalado:eq:69}
\end{equation}
Le membre de droite tend vers zéro, ce qui démontre \eqref{hmtfg:escalado:eq:68}.

\subsubsection{Enchaînement et portée de la clôture}
\label{hmtfg:escalado:17.4}

La composition effective est : lecteur dodécaphasique APP–TRIT–TPK ; profil \eqref{hmtfg:escalado:eq:1} ; transport ordonné des retours \eqref{hmtfg:escalado:eq:38}–\eqref{hmtfg:escalado:eq:40} ; classification et existence des premières racines \eqref{hmtfg:escalado:eq:50} ; localisation \eqref{hmtfg:escalado:eq:52} ; isolement par première sortie \eqref{hmtfg:escalado:eq:63} ; identité des centres à partir d’un certain rang \eqref{hmtfg:escalado:eq:66} ; changement d’échelle \eqref{hmtfg:escalado:eq:67}–\eqref{hmtfg:escalado:eq:69}. La représentation de Jacobi \eqref{hmtfg:escalado:eq:43} conserve le profil complet de ce lecteur et ses marques restent dans l’antécédent.

L’expansion spatiale du lecteur HMT et le raffinement nonadique apportent la généalogie et l’évaluation à profondeur arbitraire. Le changement d’échelle entre paramètres se démontre par la dynamique du retour, la transversalité et l’isolement que nous venons de composer. Les preuves analytiques citées interviennent sur la famille déjà construite ; leurs valeurs de comparaison ne sélectionnent ni ses phases ni ses coefficients.

La clôture concerne le sélecteur global initial du secteur uniforme \(\eta=0\), avec les dépendances théorématiques énumérées. La loi asymptotique conserve ses quantificateurs existentiels. La largeur d’une boîte de racines finies mesure la précision d’évaluation ; l’estimation \eqref{hmtfg:escalado:eq:69} exprime la convergence vers la limite.


Les programmes et certificats du supplément de reproduction permettent de reproduire les parties finies de la preuve. La continuation à toute profondeur repose sur les inductions, l’argument uniforme de première sortie et les théorèmes d’identification, avec leurs domaines explicites.
